\section{Proofs of the main results}
\label{app:proofs}
\renewcommand{\theequation}{B.\arabic{equation}}
\setcounter{equation}{0}

This appendix proves the propositions of Section~\ref{sec:population}
and the theorems of Section~\ref{sec:estimation} from the auxiliary
results of Appendix~\ref{app:aux}.

\subsection{Proof of Proposition~\ref{prop:projection}}

For fixed $t>0$, with $M$ and $K_y^\star$ as in
Lemmas~\ref{lem:app-remainder} and~\ref{lem:app-tilt} applied to
$\mathcal G$,
\[
 \frac{M(ty)}{M(y)}
 =\int_{[\underline\gamma,\Gamma]}t^{-1/g}\,K_y^\star(dg)
 \longrightarrow t^{-1/\xi}
\]
almost surely, by Lemma~\ref{lem:app-tilt} and boundedness and continuity
of $g\mapsto t^{-1/g}$ on $[\underline\gamma,\Gamma]$.
Lemma~\ref{lem:app-remainder} transfers the limit to
$\Pp(Y>ty\mid\mathcal G)/\Pp(Y>y\mid\mathcal G)$.  Convergence
simultaneously for all $t>0$ follows from rational $t$, monotonicity of
the ratio in $t$, and continuity of $t\mapsto t^{-1/\xi}$.

For the fibre identity \eqref{eq:fibre}, fix $u\in\I$ and apply the above
on the probability space $([0,1]^{p-1},K_j(u,\cdot))$ with the trivial
$\sigma$-field; $\xi_j(u)$ is by definition the essential supremum of
$v\mapsto\gamma(u,v)$ under $K_j(u,\cdot)$.  Let
$m=\max_v\gamma(u,v)$, attained by compactness and continuity.  The
essential supremum is at most $m$.  Conversely, for $\eta>0$ the set
$\{v:\gamma(u,v)>m-\eta\}$ is open and nonempty, so has positive
$K_j(u,\cdot)$-probability under (S); the essential supremum is at least
$m-\eta$ for every $\eta>0$.  \qed

\subsection{Proof of Proposition~\ref{prop:geometry}}

\emph{Claim 1.}  For $j\notin\A$, $\gamma(u,v)=g\{(u,v)_{\A}\}$ does not
involve $u$, and as $v$ ranges over the fibre the active subvector ranges
over all of $[0,1]^{s}$; by \eqref{eq:fibre},
$\xi_j(u)=\max_{w}g(w)=\gmax$ for every $u$.

\emph{Claim 2.}  By \eqref{eq:fibre}, $\xi_j(u)=\gmax$ exactly when the
fibre through $u$ intersects $\M$, i.e.\ when $u\in\pi_j(\M)$.

\emph{Claim 3.}  $u\mapsto\gmax-\xi_j(u)$ is continuous
(Lemma~\ref{lem:app-envelope-cont}), nonnegative, and by Claim~2 vanishes
on $\I$ exactly on $\I\cap\pi_j(\M)$.  A continuous nonnegative function
has positive integral over $\I$ if and only if its positivity set has
positive Lebesgue measure.

\emph{Final statement.}  Write $\M^{\A}$ for the maximizer set of $g$ in
$[0,1]^{s}$; then $\M=\{u:u_{\A}\in\M^{\A}\}$ and, for $j\in\A$,
$\pi_j(\M)$ is the projection of $\M^{\A}$, a finite set when $\M^{\A}$ is
finite; its intersection with $\I$ is Lebesgue-null and Claim~3 applies.
\qed

\subsection{Proof of Proposition~\ref{prop:margin}}

Every $x$ on the fibre $x_j=u$ has
$\|x_{\A}-u^\star_{\A}\|\ge|u-u^\star_j|$, so every point of the fibre
satisfies
$\gamma(x)\le\gmax-c\,|u-u^\star_j|^\nu$; taking the fibre maximum gives
$\xi_j(u)\le\gmax-c\,|u-u^\star_j|^\nu$.  Averaging
over $u\in\I$ gives the score bound.  For $\nu=2$,
$\E(V-u_j^\star)^2=\mathrm{Var}(V)+(1/2-u_j^\star)^2
=(1-2\varepsilon)^2/12+(1/2-u_j^\star)^2$.  \qed

\subsection{Proof of Theorem~\ref{thm:score}}

We use Proposition~\ref{thm:app-score-na} rather than
Proposition~\ref{thm:app-score}: the two differ only in that the former
averages the signed local fluctuations over the coloured rank grid
instead of bounding their supremum, which is what removes the bandwidth
from the stochastic term.

Choose $\eta_n=\sqrt{\log(pn)/n}$, $\tau=\tau_n=(pn)^{-2}$, and
$t_n=2\sqrt{\log(pn)/\mathfrak n_n}$.  The four probability terms are
then each $o(1)$: the first is $2pe^{-2\log(pn)}=2p(pn)^{-2}\to0$; the
second is $2pn(pn)^{-2}\to0$; the third is $pS_ne^{-(K_n+1)/4}\to0$
because $\log(pS_n)\le C\log(pn)=o(K_n)$ by the second condition in
\eqref{eq:rate-cond} and Lemma~\ref{lem:app-kn}; the fourth is
$4pe^{-\log(pn)}=4p(pn)^{-1}\to0$.  The second condition in
\eqref{eq:rate-cond} makes $t_n\le1$ for large $n$, and
$\Gamma t_n=O(\sqrt{\log(pn)/(n\alpha)})$ by
Remark~\ref{rem:app-hgone}.

It remains to check that the bias part of $B_n(\eta_n,t_n,\tau_n)+q_n$ in
Proposition~\ref{thm:app-score} is
$o(\Delta_{\min})$.  First, $\log(pn)=o(nh^2)$ gives
$\eta_n+\delta_n=o(h)$, so $h+\eta_n+\delta_n\le C_0h$ eventually and
$2\omega_n(h+\eta_n+\delta_n,\tau_n)\le2\omega_n(C_0h,\tau_n)
=o(\Delta_{\min})$ by (E3) and monotonicity of
$\omega_n(\cdot,\tau)$.  Second,
$(1+t_n)b_n(3\alpha)\le2b_n(3\alpha)=o(\Delta_{\min})$ by (E2).  Third,
$q_n=\omega_\xi(\delta_n)+O(1/n)$: the first part is
$o(\Delta_{\min})$ by (E1), and
$1/n\le\sqrt{\log(pn)/(n\alpha)}$ since $\alpha\le1\le n\log(pn)$, so
the second part is absorbed into the stochastic rate.  Collecting terms,
outside a vanishing probability,
\[
 \max_{j\le p}|\widehat\Psi_j-\Psi_j|
 \le a_n+b_n',\qquad
 a_n=O\Bigl(\sqrt{\log(pn)/(n\alpha)}\Bigr),\quad
 b_n'=o(\Delta_{\min}),
\]
which is the assertion.  \qed

\subsection{Proof of Theorem~\ref{thm:sis}}

Under \eqref{eq:sis-cond},
$\sqrt{\log(pn)/(n\alpha)}=o(\Delta_{\min})$, so the bound of
Theorem~\ref{thm:score} is below $\Delta_{\min}/2$ for $n$ large,
outside a vanishing probability.  On the complementary event,
Lemma~\ref{lem:app-separation} places every active score strictly before
every inactive score, so $\A\subseteq\widehat\A_d$ for every $d\ge s$.  \qed
